% Fresh section. The retained report and earlier methodology are not edited.
% For integration, include this file after Section 2 and merge references.tex.
\section{Methodology}
\label{sec:fresh_methodology}
The experiment compares five ways of classifying a clothing image: Linear, MLP,
CNN, LSTM and Transformer. Each model receives the same image information, but
processes it differently. Linear and MLP use a pixel vector, CNN uses the image
grid, LSTM reads image rows, and Transformer works with small image patches.
This section describes the implemented models and the procedure used to train
and evaluate them.

\subsection{Shared data and workflow}
Fashion-MNIST contains ten classes of $28\times28$ grayscale images
\cite{fresh_fashion}. As described in Section~2, the official training set is
divided into 54,000 training images and 6,000 validation images using a
stratified split with seed~42. The official test set contains 10,000 images.
Training fits the model weights, validation chooses which saved weights to
keep, and testing measures the selected model.

All models use the preprocessing from Section~2. Pixel values are divided by
255 and standardized with the training mean and standard deviation,
approximately 0.286209 and 0.353160. The same constants are reused for validation
and testing. No images are resized, augmented or removed.

\begin{center}
\small
\begin{tabular}{c@{\quad$\longrightarrow$\quad}c@{\quad$\longrightarrow$\quad}c}
Standardized image & Model-specific input & Ten class scores\\[5pt]
\multicolumn{3}{c}{Training updates weights; validation selects weights; testing evaluates them.}
\end{tabular}
\end{center}

Table~\ref{tab:fresh_inputs} gives the input to each model. Here, $B$ is the
number of images in a batch. Each representation contains all 784 pixel values
per image. The labels are integer class indices from 0 to 9. Every model returns
a tensor of shape $(B,10)$ containing raw class scores, called logits.

\begin{table}[h]
\centering
\caption{Inputs and trainable parameter counts of the implemented models.}
\label{tab:fresh_inputs}
\begin{tabularx}{\linewidth}{l l X r}
\toprule
Model & Input shape & Arrangement & Parameters\\
\midrule
Linear & $(B,784)$ & One vector per image & 7,850\\
MLP & $(B,784)$ & One vector per image & 235,146\\
CNN & $(B,1,28,28)$ & One grayscale image grid & 421,642\\
LSTM & $(B,28,28)$ & 28 rows, each with 28 pixels & 82,186\\
Transformer & $(B,49,16)$ & 49 patches, each with 16 pixels & 71,946\\
\bottomrule
\end{tabularx}
\end{table}

\subsection{Pixel-based baselines}
\subsubsection{Linear classifier}
The Linear model assigns a weight to each pixel for each class. For a flattened
image $x\in\mathbb{R}^{784}$, its output is
\[
z=Wx+b,\qquad W\in\mathbb{R}^{10\times784},\quad b\in\mathbb{R}^{10}.
\]
The largest score determines the prediction. This model provides a small
baseline with no hidden layer. It can learn which pixel positions support a
class, but cannot learn nonlinear combinations of those pixels.

\subsubsection{Multilayer perceptron}
The MLP adds two hidden layers with 256 and 128 units. Each hidden layer is
followed by ReLU and dropout with probability 0.2. ReLU sets negative values to
zero, while dropout randomly suppresses some activations during training.
\[
784 \;\longrightarrow\; 256 \;\longrightarrow\; 128
\;\longrightarrow\; 10.
\]
The hidden layers allow the MLP to learn nonlinear relationships among pixels.
The image is still treated as a vector: nearby pixels do not receive a special
shared operation. Dropout is disabled during validation and testing.

\clearpage
\subsection{Models for spatial and sequential structure}
\subsubsection{Convolutional neural network}
The CNN keeps the two-dimensional image layout. Convolution applies the same
small filter at different image positions, allowing the network to learn local
patterns such as edges and contours \cite{fresh_cnn}. Our model has two
convolutional blocks, followed by a dense classifier.

Both convolutions use $3\times3$ kernels, stride~1 and padding~1, so they preserve
height and width. ReLU follows each convolution. A $2\times2$ max-pooling layer
then keeps the largest value in each window and halves both spatial dimensions.
The first block produces 32 feature maps and the second produces 64.

\begin{table}[h]
\centering
\caption{CNN operations and output sizes.}
\label{tab:fresh_cnn}
\begin{tabularx}{\linewidth}{X l}
\toprule
Operation & Output shape\\
\midrule
Input image & $(B,1,28,28)$\\
$3\times3$ convolution, 32 channels; ReLU & $(B,32,28,28)$\\
$2\times2$ max pooling & $(B,32,14,14)$\\
$3\times3$ convolution, 64 channels; ReLU & $(B,64,14,14)$\\
$2\times2$ max pooling & $(B,64,7,7)$\\
Flatten & $(B,3136)$\\
Dense layer, 128 units; ReLU; dropout 0.2 & $(B,128)$\\
Dense layer, 10 outputs & $(B,10)$\\
\bottomrule
\end{tabularx}
\end{table}

The final classifier combines the learned spatial features into ten scores.
Although convolution shares weights across positions, the dense layer from
3,136 features to 128 units accounts for most of this model's parameters.
The CNN therefore has the largest parameter count in this comparison.

\subsubsection{Row-sequence LSTM}
An LSTM maintains a hidden state and a memory state while reading a sequence
\cite{fresh_lstm}. Here, the sequence comes from a static image. The first step
receives the top row, the second receives the next row, and the last receives
the bottom row. Each step contains 28 pixel values.

For a standardized image $X$, row $t$ is
$r_t=X_{0,t,:}$. The batch shape $(B,28,28)$ means
\emph{batch, row steps, pixel features}. The two dimensions of size 28 have
different roles; the sequence is not a sequence of different images.

\begin{center}
\small
\begin{tabular}{c@{\ $\longrightarrow$\ }c@{\ $\longrightarrow$\ }c@{\ $\longrightarrow$\ }c}
Row 1 & Row 2 & $\cdots$ & Row 28\\
\multicolumn{4}{c}{One shared LSTM reads the rows in order.}\\[5pt]
\multicolumn{4}{c}{$h_{28}\in\mathbb{R}^{128}
\;\longrightarrow\; \text{dropout }0.2
\;\longrightarrow\; \text{dense layer}
\;\longrightarrow\; 10\text{ scores}$}
\end{tabular}
\end{center}

The implementation uses one forward LSTM layer with 128 hidden units. Its
hidden and memory states start at zero for each image. The gates control what
information is retained and updated as successive rows are read. After the
28th row, the final hidden state is passed through dropout and a linear
classifier. No state is carried from one image to another.

This design tests whether a compact summary built from ordered rows is useful
for classification. It introduces a top-to-bottom reading order, whereas the
CNN processes local neighborhoods throughout the image.

\clearpage
\subsection{Patch-based Transformer}
The Transformer uses attention to combine information from different parts of
an image \cite{fresh_attention}. The input follows the image-patch idea
\cite{fresh_vit}, using a small configuration suited to the $28\times28$ images.

\subsubsection{Patch tokens and position information}
Each image is divided into a $7\times7$ grid of non-overlapping $4\times4$
patches. Each patch is flattened into 16 values. Patches are ordered from left
to right across each row of the grid, then from top to bottom.
A learned linear layer maps each patch from 16 values to 64 features.
A learned positional vector is added so the model can distinguish patch
locations:
\[
e_j=Wp_j+b+q_j,\qquad j=1,\ldots,49,
\]
where $p_j\in\mathbb{R}^{16}$ is a patch,
$W\in\mathbb{R}^{64\times16}$ is the projection, and
$q_j\in\mathbb{R}^{64}$ is its position vector.
Dropout with probability 0.2 is applied after this addition.

\subsubsection{Attention blocks and classification}
Two encoder blocks process the 49 tokens. Each block uses four attention heads,
a feature width of 64, and a feedforward network with hidden width 128 and
GELU activation. Layer normalization is applied before each attention and
feedforward sublayer, with residual connections around both. Dropout is 0.2.

For one attention head, the operation is
\[
\operatorname{Attention}(Q,K,V)
=\operatorname{softmax}\left(\frac{QK^\top}{\sqrt{16}}\right)V.
\]
The queries $Q$, keys $K$ and values $V$ are learned projections of the tokens.
Each head has 16 features because the 64 features are divided among four heads.
Attention weights determine how much information one patch receives from the
others. All 49 patches are available to every patch; no causal mask is used.

\begin{table}[h]
\centering
\caption{Transformer representation through the network.}
\label{tab:fresh_transformer}
\begin{tabularx}{\linewidth}{X l}
\toprule
Stage & Shape\\
\midrule
Flattened image patches & $(B,49,16)$\\
Patch projection and learned positions & $(B,49,64)$\\
Two encoder blocks & $(B,49,64)$\\
Final layer normalization; mean over patches & $(B,64)$\\
Dropout 0.2; dense classifier & $(B,10)$\\
\bottomrule
\end{tabularx}
\end{table}

The final representation is the mean of the 49 normalized output tokens.
There is no extra classification token. The encoder internally uses
sequence-first tensors of shape $(49,B,64)$; this transpose changes the layout,
not the patch order. All weights are trained from scratch.

\subsection{Learning objective}
Every model learns to assign a high score to the correct class. For image $i$,
softmax converts the logits $z_{i,k}$ into class probabilities:
\[
p_{i,k}=\frac{\exp(z_{i,k})}{\sum_{j=0}^{9}\exp(z_{i,j})},
\qquad
\mathcal{L}=-\frac{1}{B}\sum_{i=1}^{B}\log p_{i,y_i}.
\]
Here, $y_i$ is the correct class and $\mathcal{L}$ is the average cross-entropy
loss. The implementation passes logits directly to
\texttt{CrossEntropyLoss}, which performs the required calculation.
Prediction uses $\hat y_i=\arg\max_k z_{i,k}$. Class weights are not used because
the training split contains the same number of examples from each class.

\clearpage
\subsection{Training procedure}
All five models use Adam \cite{fresh_adam} with learning rate 0.001, batch size
128 and a budget of ten epochs. These settings give each model the same number
of passes over the training data. Table~\ref{tab:fresh_training} records the
shared choices.

\begin{table}[h]
\centering
\caption{Training and checkpoint selection settings.}
\label{tab:fresh_training}
\begin{tabularx}{\linewidth}{l X}
\toprule
Setting & Value\\
\midrule
Random seed & 42 for Python, NumPy, PyTorch and the training shuffle\\
Optimizer & Adam, learning rate 0.001\\
Batch size & 128; the smaller last batch is kept\\
Training duration & 10 complete epochs\\
Loss & Unweighted cross-entropy\\
Regularization & Dropout 0.2 in MLP, CNN, LSTM and Transformer\\
Checkpoint selection & Lowest validation loss; an exact tie keeps earlier weights\\
Data loading & Training shuffled; validation and test ordered; zero workers\\
Other settings & No augmentation, weight decay, gradient clipping or scheduler\\
\bottomrule
\end{tabularx}
\end{table}

Each model is initialized separately after resetting the seed. A training epoch
visits all 54,000 training images. For each batch, the program clears the
previous gradients, calculates the logits and loss, backpropagates, and updates
the weights with Adam.

At the end of the epoch, the model is evaluated on all 6,000 validation images.
Evaluation mode disables dropout, and gradient recording is turned off.
The program saves the weights whenever validation loss improves. Training
continues through epoch ten, even if the best checkpoint was found earlier.
The saved checkpoint with the lowest validation loss is then reloaded for
testing.

Epoch loss is weighted by the number of images in each batch, so a smaller final
batch does not receive the same weight as a full batch. Accuracy and macro-F1
are calculated from predictions across the complete split. Training metrics
are collected while the weights change; validation metrics use the fixed
weights at the end of an epoch.

\subsection{Evaluation measures}
The selected checkpoint is evaluated on all 10,000 official test images.
Accuracy measures the fraction classified correctly:
\[
\operatorname{Accuracy}=\frac{1}{N}\sum_{i=1}^{N}
\mathbf{1}(\hat y_i=y_i).
\]
Here, $N$ is the number of evaluated images and $\mathbf{1}$ is one when the
prediction is correct, otherwise zero.

Macro-F1 gives equal weight to every class. For class $k$, true positives
$TP_k$ are correct predictions of that class, false positives $FP_k$ are other
classes predicted as $k$, and false negatives $FN_k$ are examples of $k$
predicted as another class:
\[
F1_k=\frac{2TP_k}{2TP_k+FP_k+FN_k},
\qquad
\operatorname{MacroF1}=\frac{1}{10}\sum_{k=0}^{9}F1_k.
\]
Undefined class scores are assigned zero. Both accuracy and macro-F1 are
reported because overall accuracy alone does not show how performance varies
among categories.

The evaluation also saves cross-entropy loss, per-class precision and recall,
and a confusion matrix. Matrix rows represent true labels and columns represent
predicted labels. Correct and incorrect example images help explain recurring
mistakes. These outputs support the results and error analysis; they do not
select the training checkpoint.

\clearpage
\subsection{Runtime, reproducibility and comparison scope}
\subsubsection{Runtime measurements}
Training time covers the ten-epoch training and validation loop, including
saving improved checkpoints. It excludes data preparation before the loop and
the final checkpoint reload and validation after it.

Forward-pass timing uses a fixed batch of 128 images already on the device.
After ten warm-up calls, the program times 100 forward calls and averages their
duration. Dividing the batch time by 128 gives time per image at that batch
size. Device synchronization is performed around the timed region. This
measurement excludes data loading and host-to-device transfers and describes
batched throughput, not the latency of a single-image request.

A separate timer measures the complete test pass, including loading, transfers,
model execution, loss and prediction collection. It excludes later metric
summaries and plots. The two measurements are reported separately because they
cover different work.

\subsubsection{Reproducibility}
The implementation separates data preparation in \texttt{a1/data.py}, model
definitions in \texttt{a1/models.py}, training in \texttt{a1/train.py}, and testing
in \texttt{a1/evaluate.py}. Configuration files specify the shared training
settings. The split cache records the training and validation indices and the
normalization constants.

Each run saves the selected checkpoint, epoch history and a summary of its
configuration and environment. Evaluation checks the recorded data and split
hashes and normalization values before using the checkpoint. It also saves
per-image predictions so that the reported scores can be checked again.

The following commands illustrate the CNN workflow from the \texttt{A1} folder
with the project dependencies installed. The other models use the same entry
points and their corresponding configuration files.
\begin{quote}\small\ttfamily
python -m a1.train --model cnn --config configs/cnn.json\\
python -m a1.evaluate --checkpoint runs/m2/cnn/best.pt --device cuda
\end{quote}

\subsubsection{What the comparison can establish}
The split, preprocessing, seed, optimizer, learning rate, batch size and epoch
budget are shared. The architecture, input arrangement and parameter count
change. The experiment therefore compares these five implemented configurations.
It does not hold model capacity or computation constant.

All computational-cost measurements in Section~\ref{sec:results} use CUDA on
an NVIDIA GeForce GTX 1660 SUPER. Linear and MLP were repeated on CUDA with
the same seed, split and training protocol, in separate run directories.
The classification scores and error analysis retain the original selected
checkpoints; the baseline cost rows describe these separate CUDA repetitions.
Finally, each model has one recorded
training seed. Differences in the resulting scores describe these runs; repeated
seeds would be needed to measure how much the results vary.

